%% LyX 2.4.4 created this file.  For more info, see https://www.lyx.org/.
%% Do not edit unless you really know what you are doing.
\documentclass[oneside,english]{book}
\usepackage[T1]{fontenc}
\usepackage[latin9]{inputenc}
\setcounter{secnumdepth}{3}
\setcounter{tocdepth}{3}
\usepackage{units}
\usepackage{amstext}
\usepackage{amsthm}
\usepackage{amssymb}
\usepackage{cancel}
\usepackage{geometry}
\geometry{verbose,tmargin=1cm,bmargin=1cm,lmargin=1.5cm,rmargin=1.5cm}
\usepackage{wasysym}
\PassOptionsToPackage{normalem}{ulem}
\usepackage{ulem}

\makeatletter

%%%%%%%%%%%%%%%%%%%%%%%%%%%%%% LyX specific LaTeX commands.
%% Because html converters don't know tabularnewline
\providecommand{\tabularnewline}{\\}

%%%%%%%%%%%%%%%%%%%%%%%%%%%%%% Textclass specific LaTeX commands.
\theoremstyle{definition}
\newtheorem{example}{\protect\examplename}
\theoremstyle{definition}
\newtheorem{defn}{\protect\definitionname}
\ifx\proof\undefined\
  \newenvironment{proof}[1][\proofname]{\par
    \normalfont\topsep6\p@\@plus6\p@\relax
    \trivlist
    \itemindent\parindent
    \item[\hskip\labelsep
          \scshape
      #1]\ignorespaces
  }{%
    \endtrivlist\@endpefalse
  }
  \providecommand{\proofname}{Proof}
\fi
\theoremstyle{definition}
\newtheorem{xca}{\protect\exercisename}
\theoremstyle{plain}
\newtheorem{prop}{\protect\propositionname}
\theoremstyle{plain}
\newtheorem{thm}{\protect\theoremname}

%%%%%%%%%%%%%%%%%%%%%%%%%%%%%% User specified LaTeX commands.
\usepackage{hyperref}
\usepackage[usenames,dvipsnames]{xcolor} %adds additional color options
\hypersetup{colorlinks=true, urlcolor=Blue}%Blue

\usepackage{multicol}

\usepackage{tikz}
\usepackage{tikz-3dplot}
\usetikzlibrary{calc,arrows,shapes,backgrounds,matrix,shapes.geometric,fit,trees,positioning}

\usetikzlibrary{patterns}
\usetikzlibrary{plotmarks}
\usepackage{pgfplots}
\pgfplotsset{compat=newest}

\usepackage{calc}
\usepackage[nomessages]{fp}

\usepackage{datetime}

%%%%%commands
\newcommand{\mb}[1]{$\mathbb{#1}$}
\newcommand{\funct}[2]{$f\colon#1\rightarrow#2$}
% Added by lyx2lyx
\setlength{\parskip}{\smallskipamount}
\setlength{\parindent}{0pt}

\makeatother

\usepackage{babel}
\providecommand{\definitionname}{Definition}
\providecommand{\examplename}{Example}
\providecommand{\exercisename}{Exercise}
\providecommand{\propositionname}{Proposition}
\providecommand{\theoremname}{Theorem}

\begin{document}

\chapter{Functions and Relations}

\section{Functions}

The previous chapter covered sets and their basic operations. Now
we introduce the notion of a function between sets. Just as in calculus,
a function will assign a unique output for every input. Familiar notations
such as, $f$ and $f(x)$ will be used. The same concepts will be
used; only the inputs and outputs will include things other than numbers.
A \emph{function }(or \emph{mapping) }from sets $A$ to $B$\emph{,
$f\colon A\rightarrow B$,} is a rule that assigns each element in
$A$ to exactly one element in $B$.
\begin{itemize}
\item $A$ is called the \emph{domain }of $f$. 

\begin{itemize}
\item $B$ is called the \emph{range }of $f$.\footnote{Also called the \emph{co-domain.}}
\item If $f(a)=b$ for some $a\in A$, then $b$ is called the\emph{ image
of $a$ under $f$.}
\item $\{b\colon b\in B\mbox{ and }f(a)=b\mbox{ for some }a\in A\}$ is
called the \emph{image of} $A$ \emph{under $f$}.
\end{itemize}
\item Functions are well defined. Meaning that an input is assigned to one
and only output. 
\end{itemize}
\begin{example}
Let $f\colon\mathbb{R}\rightarrow\mathbb{R}$. Then $f(x)=x^{2}+1$
is a function. it is clear from the equation that $f$ is well defined.
\end{example}

\begin{example}
Let $f\colon\mathbb{R}\rightarrow\mathbb{R}$. Then $f(x)=\pm\sqrt{x}$
is \uline{not} a function. Since $f(1)=+1$ and $-1$.
\end{example}

\begin{example}
\label{ex-floor-funct}Consider $f\colon\mathbb{R}\rightarrow\mathbb{Z}$
given by the \emph{floor function $f(x)=\left\lfloor x\right\rfloor $
}where $\left\lfloor x\right\rfloor $ is the largest integer less
than or equal to $x$, e.g., $\left\lfloor \pi\right\rfloor =3$.
Note that the floor function is also a function from $\mathbb{R}$
to $\mathbb{R}$, $f\colon\mathbb{R}\rightarrow\mathbb{Z}$.
\end{example}

\begin{example}
Let $f\colon\mathbb{Z}\times\mathbb{Z}\rightarrow\mathbb{Z}$ given
by $f(x,y)\rightarrow xy$. Then $f$ is a function. 
\end{example}

\begin{example}
Let $A=\{a,b,c,\dots,z\}$ and $f\colon A\rightarrow\mathbb{Z}$ be
function defined by assigning the $n^{\mbox{th}}$letter to $n$,
e.g., $f(3)=c$.
\end{example}
\begin{defn}
Let $f\colon B\rightarrow C$ and $g\colon A\rightarrow B$ be functions,
then the \emph{composition }of $f$ and $g$ is $f\circ g\colon A\rightarrow C$
defined as $f(g(a))=c$ where $g(a)=b$ and $f(b)=c$ for elements
$a\in A$, $b\in B$, and $c\in C$.
\end{defn}
\begin{figure}
\begin{centering}
\begin{tikzpicture}[description/.style={fill=white,inner sep=2pt}]   
	\matrix (m)[matrix of math nodes,     
	row sep=3em,column sep=2.5em,     
	text height=1.5ex, text depth=0.25ex,     	nodes={ellipse,draw,minimum width=2cm},   	]
	{
		A && B \\     
		& C & 
		\\   };   
	\path[->,font=\scriptsize]   (m-1-1) edge node[auto] {$ f $} (m-1-3)   edge node[description] {$ g \circ f $} (m-2-2)   (m-1-3) edge node[auto] {$ g $} (m-2-2); \end{tikzpicture}
\par\end{centering}
\caption{Composition of functions $f$ and $g$.}
\end{figure}

\begin{itemize}
\item In general the composition of function is not commutative, i.e., $f\circ g\neq g\circ f$.
\end{itemize}
\begin{example}
Let $f,g:\mathbb{R}\rightarrow\mathbb{R}$ and $f(x)=2x+1$ and $g(x)=x^{2}$.
Then $f\circ g(x)=f(g(x))=2(x^{2})+1=2x^{2}+1$.
\end{example}

\begin{defn}
A function, \funct{A}{B}, is called \emph{one-to-one}\footnote{Also called an \emph{injective function}}
if $f(a_{1})=f(a_{2})$ implies that $a_{1}=a_{2}$.
\end{defn}
\begin{example}
$f:\mathbb{Z}\rightarrow\mathbb{Z}$ where $f(x)=x^{3}$ is one-to-one.
\end{example}
\begin{proof}
Suppose that $f(x_{1})=f(x_{2})$ for some $x_{1},x_{2}\in\mathbb{Z}$.
So $(x_{1})^{3}=(x_{2})^{3}\Rightarrow\left((x_{1})^{3}\right)^{\nicefrac{1}{3}}=\left((x_{2})^{3}\right)^{\nicefrac{1}{3}}\Rightarrow x_{1}=x_{2}$
\end{proof}

\begin{xca}
Show that any linear function $f:\mathbb{R}\rightarrow\mathbb{R}$
with a non-zero defined slope is one-to-one. 
\end{xca}
\begin{example}
$f:\mathbb{R}\rightarrow\mathbb{R}$ where $f(x)=|x|$ is not one-to-one.
\end{example}
\begin{itemize}
\item []$|-1|=1=|1|$ but $-1\neq1$ so $f$ is not one-to-one.
\end{itemize}
\begin{prop}
\label{prop:f is 1-1 then f^-1 is fuct}If $f\colon A\rightarrow B$
then $f^{-1}=\{x\colon x\in A\mbox{, }f(x)=b\}$ is a function if
and only if $f$ is one-to-one.
\end{prop}
\begin{xca}
Prove proposition \ref{prop:f is 1-1 then f^-1 is fuct}.
\end{xca}
\begin{defn}
Let $f\colon A\rightarrow B$ be a one-to-one function. We call the
function guarteed to exist by proposition \ref{prop:f is 1-1 then f^-1 is fuct},
$f^{-1}\colon B\rightarrow A$, the \emph{inverse} of $f$.
\end{defn}

\begin{defn}
A function \funct{A}{B} is \emph{onto }if for every $b\in B$ there
is an $a\in A$ such that $f(a)=b$.
\end{defn}
\begin{example}
$f\colon\mathbb{R}\rightarrow\mathbb{R}$ where $f(x)=x^{2}$ is onto.
\end{example}
\begin{proof}
If $b\in\mathbb{R}$ then $\sqrt{b}\in\mathbb{R}$, and $f(\sqrt{b})=\left(\sqrt{b}\right)^{2}=b$.
\end{proof}

\begin{example}
$f\colon\mathbb{Z}\rightarrow\mathbb{R}$ where $f(x)=x^{2}$ is not
onto.
\end{example}
\begin{proof}
$2\in\mathbb{R}$ but $\sqrt{2}\notin\mathbb{Z}$.
\end{proof}
\begin{xca}
Prove that the floor function (see example \ref{ex-floor-funct})
is onto but not one-to-one. 
\end{xca}
\begin{figure}[h]
\hfill%
\begin{minipage}[t]{0.33\textwidth}%
 \begin{tikzpicture}[scale=.6,ele/.style={fill=black,circle,minimum width=.8pt,inner sep=1pt},every fit/.style={ellipse,draw,inner sep=-2pt}]   
	\node[label=center:$\textbf{A}$] (a0) at (0,5) {};
	\node[ele,label=left:$a$] (a1) at (0,4) {};
	\node[ele,label=left:$a$] (a1) at (0,4) {};       
	\node[ele,label=left:$b$] (a2) at (0,3) {};       
	\node[ele,label=left:$c$] (a3) at (0,2) {};   
	\node[ele,label=left:$d$] (a4) at (0,1) {};
	\node[label=center:$\textbf{B}$] (b0) at (4,5) {};	
	\node[ele,,label=right:$1$] (b1) at (4,4) {};   
	\node[ele,,label=right:$2$] (b2) at (4,3) {};   
	\node[ele,,label=right:$3$] (b3) at (4,2) {};   
	\node[ele,,label=right:$4$] (b4) at (4,1) {};
	\node[draw,fit= (a1) (a2) (a3) (a4),minimum width=2cm] {};
	\node[label=center:$f$] 
(f1) at (2,4) {};
	 
	\node[draw,fit= (b1) (b2) (b3) (b4),minimum width=2cm] {};     
	
	\draw[->,thick,shorten <=2pt,shorten >=2pt] (a1) -- (b2);   
	\draw[->,thick,shorten <=2pt,shorten >=2] (a2) -- (b1);   
	\draw[->,thick,shorten <=2pt,shorten >=2] (a3) -- (b3);   
	\draw[->,thick,shorten <=2pt,shorten >=2] (a4) -- (b4);  
\end{tikzpicture}\caption{}{\funct{A}{B} is  one-to-one and onto.}%
\end{minipage}\hfill%
\begin{minipage}[t]{0.33\textwidth}%
\begin{center}
 \begin{tikzpicture}[scale=.6,ele/.style={fill=black,circle,minimum width=.8pt,inner sep=1pt},every fit/.style={ellipse,draw,inner sep=-2pt}]   
	\node[label=center:$\textbf{A}$] (a0) at (0,5) {};
	\node[ele,label=left:$a$] (a1) at (0,4) {};
	\node[ele,label=left:$a$] (a1) at (0,4) {};       
	\node[ele,label=left:$b$] (a2) at (0,3) {};       
	\node[ele,label=left:$c$] (a3) at (0,2) {};   
	\node[ele,label=left:$d$] (a4) at (0,1) {};
	\node[label=center:$\textbf{B}$] (b0) at (4,5) {};	
	\node[ele,,label=right:$1$] (b1) at (4,4) {};   
	\node[ele,,label=right:$2$] (b2) at (4,3) {};   
	\node[ele,,label=right:$3$] (b3) at (4,2) {};   
	\node[ele,,label=right:$4$] (b4) at (4,1) {};
	\node[draw,fit= (a1) (a2) (a3) (a4),minimum width=2cm] {};
	\node[label=center:$f$] 
(f1) at (2,4) {};
	 
	\node[draw,fit= (b1) (b2) (b3) (b4),minimum width=2cm] {};     
	
	\draw[->,thick,shorten <=2pt,shorten >=2pt] (a1) -- (b2);   
	\draw[->,thick,shorten <=2pt,shorten >=2] (a2) -- (b2);   
	\draw[->,thick,shorten <=2pt,shorten >=2] (a3) -- (b4);   
	\draw[->,thick,shorten <=2pt,shorten >=2] (a4) -- (b3);  
\end{tikzpicture}\caption{}{\funct{A}{B} is not one-to-one and is not onto.}
\par\end{center}%
\end{minipage}\hfill%
\begin{minipage}[t]{0.33\textwidth}%
\begin{center}
 \begin{tikzpicture}[scale=.6,ele/.style={fill=black,circle,minimum width=.8pt,inner sep=1pt},every fit/.style={ellipse,draw,inner sep=-2pt}]   
	\node[label=center:$\textbf{A}$] (a0) at (0,5) {};
	\node[ele,label=left:$a$] (a1) at (0,4) {};
	\node[ele,label=left:$a$] (a1) at (0,4) {};       
	\node[ele,label=left:$b$] (a2) at (0,3) {};       
	\node[ele,label=left:$c$] (a3) at (0,2) {};   
	\node[ele,label=left:$d$] (a4) at (0,1) {};
	\node[label=center:$\textbf{B}$] (b0) at (4,5) {};	
	%\node[ele,,label=right:$1$] (b1) at (4,4) {};   
	\node[ele,,label=right:$2$] (b2) at (4,3) {};   
	\node[ele,,label=right:$3$] (b3) at (4,2) {};   
	\node[ele,,label=right:$4$] (b4) at (4,1) {};
	\node[draw,fit= (a1) (a2) (a3) (a4),minimum width=2cm] {};
	\node[label=center:$f$] 
(f1) at (2,4) {};
	 
	\node[draw,fit= (b1) (b2) (b3) (b4),minimum width=2cm] {};     
	
	\draw[->,thick,shorten <=2pt,shorten >=2pt] (a1) -- (b2);   
	\draw[->,thick,shorten <=2pt,shorten >=2] (a2) -- (b2);   
	\draw[->,thick,shorten <=2pt,shorten >=2] (a3) -- (b4);   
	\draw[->,thick,shorten <=2pt,shorten >=2] (a4) -- (b3);  
\end{tikzpicture}\caption{}{\funct{A}{B} is onto but not one-to-one.}
\par\end{center}%
\end{minipage}\hfill
\end{figure}

\begin{defn}
A function that is both one-to-one and onto is called a \emph{bijection.}
\end{defn}
\begin{example}
The function \funct{\mathbb{Z}}{\mathbb{Z}} defined by $f(x)=x+1$
is a bijection.

\begin{proof}
Let $b\in\mathbb{Z}$ then $f(b-1)=b-1+1=b$ so $f$ is onto. Now
suppose that $f(a_{1})=f(a_{2})$. $f(a_{1})=a_{1}+1=a_{2}+1=f(a_{2})$
implies that $a_{1}=a_{2}$. So then $f$ is one-to-one. 
\end{proof}
\end{example}
\begin{thm}
\textbf{\textup{Pigeonhole Principle.}}\textbf{\textup{\emph{ }}}\textup{\emph{If
$n>m$ and $n$ items are placed in $m$ boxes then at least one box
will have two items.}}

Equivalently: If $f\colon A\rightarrow B$ and $|A|>|B|$ then $f$
is not one-to-one.
\end{thm}
\begin{itemize}
\item The pigeonhole principle says that there is at least one element $b\in B$
such that $f(a)=f(a')=b$ for some $a,a'\in A$.

\begin{proof}
Suppose to the contrary that $|A|=n>m=|B|$ and there is a a one-to-one
function $f$ from $A$ to $B$. If $A=\{a_{1},a_{2},\dots,a_{n}\}$
then $f(A)=\{f(a_{1}),f(a_{2}),\dots,f(a_{n})\}$ lists $n$ distinct
elements since by definition $f(a_{j})\neq f(a_{k})$. Then $\{f(a_{1}),f(a_{2}),\dots,f(a_{n})\}=\{b_{1},b_{2},\dots,b_{n}\}$
lists $n$ distinct elements in $B$, but this would mean $n\leq m$
which is a contradiction.
\end{proof}
\end{itemize}
\begin{thm}
\label{thm:order and 1-1 and onto}Let $f$ be a function from $A$
to $B$. The the following hold:

\begin{enumerate}
\item If $f$ is one-to-one then $|B|\geq|A|$.
\item If $f$ is onto then $|A|\geq|B|$.
\item If $f$ is a bijection then $|A|=|B|$.
\end{enumerate}
\end{thm}
Part 1 is true by the pigeonhole principle. Part 2 and 3 are left
as exercises.
\begin{example}
Prove parts 2 and 3 of theorem \ref{thm:order and 1-1 and onto}.
\end{example}
\begin{xca}
Let $f\colon A\rightarrow B$. Prove that if $f$ is onto then $|A|\geq|B|$.
\end{xca}

\begin{xca}
Let $f\colon A\rightarrow B$. Prove that if $f$ is onto then $|A|\geq|B|$.
\end{xca}

\begin{xca}
$\bigstar$ Let the function $f:A\rightarrow B$ be 1-1 and onto,
i.e., $f^{-1}$ exists and $\left|A\right|=\left|B\right|$. Prove
that $\left(f^{-1}\circ f\right):A\rightarrow A$ is the identity
function on $A$ and that $\left(f\circ f^{-1}\right):B\rightarrow B$
is the identity function on $B$.
\end{xca}

\begin{xca}
$\bigstar$ Let $f:A\rightarrow B$ be 1-1 and onto. Prove that $g=f^{-1}$,
if $\left(g\circ f\right):A\rightarrow A$ is the identity function
on $A$ and $\left(f\circ g\right):B\rightarrow B$ is the identity
function on $B$. 
\end{xca}

\begin{xca}
[$\smiley$\textbf{.}]Two sets are said to be \emph{equal }if there
exists a bijective function between them; we write $\left|A\right|=\left|B\right|$.
$B$ is said to have \emph{strictly greater cardinality} than $B$
if there is a one-to-one function from $A$ into $B$, but there does
not exist a bijective function from $A$ into $B$; we write$\left|A\right|<\left|B\right|$.
Prove that $\left|\mathbb{N}\right|<\left|\mathbb{R}\right|$.
\end{xca}

\section{Relations}

In a relation, things are either related or they are not. For example,
the statement ``the woman $x$ is taller than the man $y$'' is
either true or false for any $x$ from the set of woman and $y$ for
the set of men. This statement is either true or false. Relations
are a generalization of the arithmetic relations such as $=$, $>$,
$\leq$, etc.

It is different than the concept of function as it does not determine
particular parings of elements. In the example there may be many woman
taller than any given man. We will investigate binary relations, meaning
we will apply statements to two elements at a time. 
\begin{defn}
Let $A$ and $B$ be sets and $P(a,b)$ be true or false for any $a\in A$
and $b\in B$. Then $\mathcal{R}(A,B,P(a,b))$ is called a \emph{relation.
$P(a,b)$ }is called a \emph{statement}. 
\end{defn}
\begin{itemize}
\item If $P(a,b)$ is true we write $a\mathcal{R}b$ and say that \emph{a
is related to b}.
\item If $P(a,b)$ is false we write $a\cancel{\mathcal{R}}b$ and say that
\emph{a is not related to b.}
\item When $A$, $B$, and $P(a,b)$ are understood, we will write just
$\mathcal{R}$.
\end{itemize}
\begin{example}
Let $A$ be the set of men, $B$ be the set of women, and $P(a,b)$
be the statement ``$x$ is $y$'s husband''

\begin{itemize}
\item []Then Mr.\,Obama$\mathcal{R}$Mrs.\,Obama
\item []Then Mr.\,Obama$\not\mathcal{R}$Mr.\,Rommney
\end{itemize}
\end{example}

\begin{example}
Let $V=\{1,2,3,4\}$ and $U=\{a,b,c\}$ and $\mathcal{R}=\{(1,a),(1,b),(4,c)\}$.

\begin{itemize}
\item []$1\mathcal{R}a$, $1\mathcal{R}b$, and $4\mathcal{R}c$
\item []$1\cancel{\mathcal{R}}c$, $2\cancel{\mathcal{R}}a$, and $3\cancel{\mathcal{R}}c$
(among many)
\end{itemize}
\end{example}
In the above example we did not even need to know the statement; $\mathcal{R}$
told us the relations. We can think of $\mathcal{R}$ as a set of
pairs, but note that it is not a function from $V$ to $U$. Since
then it would map $1\rightarrow a$ and $1\rightarrow b$ -it is not
well defined. 
\begin{xca}
Let $A=\{1,2,3,4\}$, $B=\{5,6\}$, and define $P(a,b)=\mbox{''\emph{a }divides }b\mbox{"}$
for $a\in A$ and $b\in B$. Find $\mathcal{R}(A,B,P)$.
\end{xca}

\begin{example}
Let $A=\{1,2,3\}$ and define $P(a,b)=\mbox{''}a<b\mbox{"}$. Then
$\mathcal{R}=\{(1,2),(1,3),(2,3)\}$ since $1<2$, $1<3$, and $2<3$
are all true.
\end{example}
\begin{itemize}
\item If the two sets in binary relation $\mathcal{R}$ are the same, $\mathcal{R}(A,A,P)$,
then we say that $\mathcal{R}$ \emph{is a relation over $A$.}\footnote{A binary relation over a single set is also called an \emph{endorelation}.}
\end{itemize}
\begin{defn}
Let $\mathcal{R}$ be a relation over a set $A$ then:

\begin{enumerate}
\item If $a\mathcal{R}a$ for all $a\in A$ then $\mathcal{R}$ is \emph{reflexive}.

\begin{enumerate}
\item []equivalently: $(a,a)\in\mathcal{R}$ for all $a\in\mathcal{R}\Rightarrow$$\mathcal{R}$
is \emph{reflexive}.
\end{enumerate}
\item If $a\mathcal{R}b$ implies that $b\mathcal{R}a$ for all $a,b\in A$
then $\mathcal{R}$ is \emph{symmetric}.

\begin{enumerate}
\item []equivalently: $(a,b)\in\mathcal{R}$$\Rightarrow(b,a)\in\mathcal{R}$
for all $a,b\in\mathcal{R}\Rightarrow$$\mathcal{R}$ is \emph{symmetric}.
\end{enumerate}
\item If $a\mathcal{R}c$ whenever $a\mathcal{R}b$ and $b\mathcal{R}c$
for all $a,b,c\in A$ then $\mathcal{R}$ is \emph{transitive}.

\begin{enumerate}
\item []equivalently: $(a,c)\in\mathcal{R}$ whenever $(a,b)\in\mathcal{R}$
and $(b,c)\in\mathcal{R}$ or all $a,b,c\in\mathcal{R}\Rightarrow$$\mathcal{R}$
is \emph{transitive}.
\end{enumerate}
\end{enumerate}
\end{defn}
\begin{example}
Let $\mathcal{R}="|"$ over $\mathbb{Z}$; that is, let $\mathcal{R}$
be defined by the statement ``$x$ divides $y$'' where $x,y\in\mathbb{Z}$.

\begin{enumerate}
\item Is $\mathcal{R}$ reflexive ? Yes. Let $a\in\mathbb{Z}$ if $a|a$
then $a|a$.
\item Is $\mathcal{R}$ symmetric? No. $2|4$ but $4\cancel{|}2$. 
\item Is $\mathcal{R}$ transitive? Yes. Let $a,b,c\in\mathbb{Z}$ such
that if $a|b$ and $b|c$. We can then write $c=bn$ and $b=am$ for
some $m,n\in\mathbb{Z}$. If $b=am$ then $c=(am)n$ which implies
that $a|c$.
\end{enumerate}
\end{example}

\begin{defn}
An \emph{equivalence relation} is a relation which is reflexive, symmetric,
and transitive.
\end{defn}
\begin{example}
Let $\mathcal{R}$ be ``equals'' over $\mathbb{R}$. That is $(a,b)\in\mathcal{R}$
if $a=b$. $\mathcal{R}$ is an equivalence relation. 

\begin{itemize}
\item []Reflexive: $a=a\Rightarrow a=a$.
\item []Symmetric: $a=b\Rightarrow b=a$.
\item []Transitive: suppose that $a=b$ and $b=c$ then $a=b=c\Rightarrow a=c$.
\end{itemize}
\end{example}
\begin{xca}
Each of the following statements, $P(x,y)$, define a relation $R$
in the natural numbers, $\mathbb{N}$, i.e., $\mathcal{R}=(\mathbb{N},\mbox{\ensuremath{\mathbb{N}}},P(x,y))$.
Show whether each of the following relations is reflexive. 

\begin{enumerate}
\item $P(x,y)=$$x$ is less than or equal to $y$
\item $P(x,y)=$$x$ divides $y$
\item $x+y=10$
\item $x$ and $y$ are relatively prime.
\end{enumerate}
\end{xca}

\begin{table}
\noindent Use the following to answer the next two problems. Let $E=\{1,2,3\}$.
Consider the following relations in $E$.

\noindent{}%
\begin{tabular}{ll}
$\mathcal{R}_{1}=\left\{ \left(1,1\right),\left(2,1\right),\left(2,2\right),\left(3,2\right),\left(2,3)\right)\right\} $ & $\mathcal{R}_{4}=\left\{ \left(1,1\right),\left(3,2\right),\left(2,3\right)\right\} $\tabularnewline
$\mathcal{R}_{2}=\left\{ \left(1,1\right)\right\} $ & $\mathcal{R}_{5}=E\times E$\tabularnewline
$\mathcal{R}_{3}=\left\{ \left(1,2\right)\right\} $ & \tabularnewline
\end{tabular}

\caption{Use the above relations to complete}
\end{table}

\begin{xca}
Let $E=\{1,2,3\}$. Consider the following relations over $E$.\\
\begin{tabular}{ll}
$\mathcal{R}_{1}=\left\{ \left(1,1\right),\left(2,1\right),\left(2,2\right),\left(3,2\right),\left(2,3)\right)\right\} $ & $\mathcal{R}_{4}=\left\{ \left(1,1\right),\left(3,2\right),\left(2,3\right)\right\} $\tabularnewline
$\mathcal{R}_{2}=\left\{ \left(1,1\right)\right\} $ & $\mathcal{R}_{5}=E\times E$\tabularnewline
$\mathcal{R}_{3}=\left\{ \left(1,2\right)\right\} $ & \tabularnewline
\end{tabular} \\
Determine whether relation $\mathcal{R}_{1},\mathcal{R}_{2},\mathcal{R}_{3},\mathcal{R}_{4},\mathcal{R}_{5}$
is symmetric.
\end{xca}

\begin{xca}
Let $E=\{1,2,3\}$. Consider the following relations in $E$. Consider
the following relations over $E$.\\
\begin{tabular}{ll}
$\mathcal{R}_{6}=\left\{ \left(1,2\right),\left(2,2\right)\right\} $ & $\mathcal{R}_{9}=\left\{ \left(1,1\right)\right\} $\tabularnewline
$\mathcal{R}_{7}=\left\{ \left(1,2\right),\left(2,3\right),\left(1,3\right),\left(2,1\right),\left(1,1\right)\right\} $ & $\mathcal{R}_{10}=E\times E$\tabularnewline
$\mathcal{R}_{8}=\left\{ \left(1,2\right)\right\} $ & \tabularnewline
\end{tabular}\\
Determine whether each relation $\mathcal{R}_{6},\mathcal{R}_{7},\mathcal{R}_{8},\mathcal{R}_{9},\mathcal{R}_{10}$
is transitive. 
\end{xca}

\begin{xca}
Define the following relation on the real numbers, $\mathcal{R}:x\mathcal{R}y\iff xy\geq0$.
Prove that $\mathcal{R}$ is reflexive and symmetric, but that $\mathcal{R}$
is not an equivalence relation.
\end{xca}

\begin{xca}
Define the following equivalence relation on the set of subsets of
$X$ as follows: $A\mathcal{R}B\iff A\cap B\neq\emptyset$ for $A,B\subseteq X$
Prove or disprove that $\mathcal{R}$ is an equivalence relation. 
\end{xca}

Let $\mathcal{R}$ and $\mathcal{R}'$ be symmetric relations in a
set A. Prove that $\mathcal{R}\cap\mathcal{R}'$ is a symmetric relation
in A. (note that $\mathcal{R}$ and $\mathcal{R}'$ are subsets of
$A\times A$).

\begin{xca}
$\bigstar$ Prove: Let $R$ and $R'$ be symmetric relations in a
set $A$; then $R\cap R'$ is a symmetric relation in $A$.
\end{xca}

\begin{xca}
$\bigstar$ Define $a\mathcal{R}b$ as $a=b\bmod n$ for $a,b\in\mathbb{Z}$
and $n\in\mathbb{Z}^{+}$.\footnote{$a=b\bmod n$ if $n|a-b$ for integers $n,a,b$. For example, $5=1\bmod4$
and $3=3\bmod4$. There is an infinite class of integers which meeting
this criteria; this set is notated $[a]=\{a+kn\colon k\in\mathbb{Z}\}$.
Modular arithemetic is covered in chapter \ref{chap:Congruence}.} $a=b\bmod n$ is an equivalence relation.
\end{xca}

\end{document}
